\documentclass[12pt]{article}

%----------------- Packages ------------------
\usepackage[utf8]{inputenc}
\usepackage[T1]{fontenc}
\usepackage{amsmath, amssymb}
\usepackage{geometry}
\usepackage{xcolor}
\usepackage{titlesec}
\usepackage{fancyhdr}
\usepackage[breakable]{tcolorbox}
\usepackage{graphicx}
\usepackage{tikz}
\usepackage{float}
\usetikzlibrary{arrows.meta, decorations.markings}

%----------------- Page Setup -----------------
\geometry{a4paper, margin=2.5cm}
\pagestyle{fancy}
\fancyhf{}
\rhead{Final Exam — \textbf{2021}}
\lhead{probability and statistics}
\cfoot{\thepage}

%----------------- Title Styling --------------
\titleformat{\section}{\normalfont\Large\bfseries}{Exercise \thesection:}{1em}{}
\titleformat{\subsection}{\normalfont\bfseries}{Answer:}{1em}{}

%----------------- Custom Boxes ----------------
\tcbuselibrary{listingsutf8}
\newtcolorbox{answerbox}{
  colback=gray!10,
  colframe=black,
  fonttitle=\bfseries,
  title=Answer Area,
  breakable,
  before skip=10pt,
  after skip=10pt
}

%----------------- Document Start --------------
\begin{document}

%----------------- Exam Info -------------------
\begin{center}
  \Large\textbf{UNIVERSITY NAME} \\[1em]
  \large\textit{probability and statistics — Final Exam} \\[0.5em]
  \large\textit{Year: 2021} \\[2em]
\end{center}

\vspace{0.5cm}

%----------------- EXERCISE 1 ------------------
\section{}
A parking lot has 10 spaces numbered from 1 to 10. The number of cars \( N \) arriving at the lot in one hour follows a Poisson distribution with parameter \( \lambda > 0 \). It is assumed that drivers choose a space at random and independently of one another.

Let \( G \) be the event “the driver chooses space number 1” and \( X \) be the random variable “number of drivers who choose space number 1”.
\\
$\text{Calculate}$:
\begin{enumerate}
    \item \( E(N) \)
    \item \( P(G) \)
    \item \( P(X = k \mid N = m) \) (distinguish the cases \( k > m \) and \( k \leq m \))
    \item \( P(X = k) \)
    \item \( E(X) \) and \( \text{Var}(X) \)
    \item The probability generating function of \( N \) and that of \( X \).
\end{enumerate}

\begin{answerbox}
\begin{enumerate}
  \item \( E(N) \)

\( N \sim \mathcal{P}(\lambda) \), so:
\[
E(N) = \lambda.
\]

  \item \( P(G) \)

  Each driver chooses a space uniformly among 10 spaces:
\[
P(G) = \frac{1}{10}.
\]

  \item \( P(X = k \mid N = m) \)

  Given \( N = m \), drivers choose independently. \( X \mid N = m \sim \text{Binomial}(m, \frac{1}{10}) \).  
Thus:
\[
P(X = k \mid N = m) = 
\begin{cases}
\binom{m}{k} \left( \frac{1}{10} \right)^k \left( \frac{9}{10} \right)^{m-k} & \text{if } 0 \le k \le m, \\
0 & \text{if } k > m.
\end{cases}
\]

  \item \( P(X = k) \)

  By total probability:
\begin{align*}
  P(X = k)  & = \sum_{m=k}^{\infty} P(X = k \mid N = m) P(N = m) \\
            & = \sum_{m=k}^{\infty} \binom{m}{k} \left( \frac{1}{10} \right)^k \left( \frac{9}{10} \right)^{m-k} \cdot \frac{e^{-\lambda} \lambda^m}{m!}.
\end{align*}
Let \( j = m - k \), then \( m = k + j \):
\[
P(X = k) = e^{-\lambda} \left( \frac{1}{10} \right)^k \sum_{j=0}^{\infty} \binom{k+j}{k} \left( \frac{9}{10} \right)^{j} \frac{\lambda^{k+j}}{(k+j)!}.
\]
Note: \(\binom{k+j}{k} \frac{1}{(k+j)!} = \frac{1}{k! \, j!}\). So:
\[
P(X = k) = e^{-\lambda} \frac{(\lambda/10)^k}{k!} \sum_{j=0}^{\infty} \frac{(9\lambda/10)^j}{j!}
= e^{-\lambda} \frac{(\lambda/10)^k}{k!} e^{9\lambda/10}.
\]
Thus:
\[
P(X = k) = \frac{e^{-\lambda/10} (\lambda/10)^k}{k!}.
\]

  \item \( E(X) \) and \( \text{Var}(X) \)

  Since \( X \sim \mathcal{P}(\lambda/10) \):
\[
E(X) = \frac{\lambda}{10}, \quad \text{Var}(X) = \frac{\lambda}{10}.
\]

  \item Probability generating functions

  For \( N \): \( G_N(s) = E(s^N) = e^{\lambda(s-1)} \).  

  For \( X \): Since \( X \sim \mathcal{P}(\lambda/10) \), \( G_X(s) = e^{(\lambda/10)(s-1)} \). 
\end{enumerate}
\end{answerbox}

%----------------- EXERCISE 2 ------------------
\section{}
Let \( X \) be a uniformly distributed random variable on the interval \([-1, 1]\).

\begin{enumerate}
    \item Find the distribution function of the random variable \( Y = 4-X^2 \).
    \item Compute \( P(2 \leq Y \leq 3.5) \).
\end{enumerate}

\begin{answerbox}
\( X \sim \mathcal{U}([-1,1]) \), so \( f_X(x) = \frac{1}{2} \mathbf{1}_{[-1,1]}(x) \).

\begin{enumerate}
  \item Distribution function of \( Y = 4 - X^2 \)

  Note: \( Y \) takes values in \([3, 4]\), since \( X^2 \in [0,1] \). For \( y \in [3,4] \):
\[
F_Y(y) = P(Y \le y) = P(4 - X^2 \le y) = P(X^2 \ge 4 - y).
\]
Let \( t = 4 - y \in [0,1] \). Then:
\[
F_Y(y) = P(X^2 \ge t) = P(|X| \ge \sqrt{t}) = 1 - P(|X| < \sqrt{t}).
\]
Since \( f_X \) is symmetric:
\[
P(|X| < \sqrt{t}) = \int_{-\sqrt{t}}^{\sqrt{t}} \frac{1}{2} \, dx = \frac{1}{2} \cdot 2\sqrt{t} = \sqrt{t}.
\]
Thus:
\[
F_Y(y) = 1 - \sqrt{4 - y}, \quad y \in [3,4].
\]
Outside: \( F_Y(y) = 0 \) for \( y < 3 \), and \( F_Y(y) = 1 \) for \( y > 4 \).

The density: for \( y \in (3,4) \):
\[
f_Y(y) = \frac{d}{dy} F_Y(y) = \frac{d}{dy} \left( 1 - (4-y)^{1/2} \right) = \frac{1}{2} (4-y)^{-1/2}.
\]
Thus:
\[
F_Y(y) = 
\begin{cases}
0 & y < 3, \\
1 - \sqrt{4-y} & 3 \le y \le 4, \\
1 & y > 4,
\end{cases}
\quad
f_Y(y) = 
\begin{cases}
\frac{1}{2\sqrt{4-y}} & 3 < y < 4, \\
0 & \text{otherwise}.
\end{cases}
\]

  \item Compute \( P(2 \le Y \le 3.5) \)

  Since \( Y \ge 3 \), this is \( P(3 \le Y \le 3.5) \). Using the CDF:
\[
P(3 \le Y \le 3.5) = F_Y(3.5) - F_Y(3) = \left(1 - \sqrt{4 - 3.5}\right) - \left(1 - \sqrt{4-3}\right).
\]
Compute: \( 4 - 3.5 = 0.5 \), \( \sqrt{0.5} = \frac{\sqrt{2}}{2} \). \( 4 - 3 = 1 \), \( \sqrt{1} = 1 \).  
So:
\[
P = \left(1 - \frac{\sqrt{2}}{2}\right) - (1 - 1) = 1 - \frac{\sqrt{2}}{2} - 0 = 1 - \frac{\sqrt{2}}{2}.
\]
\end{enumerate}
\end{answerbox}

%----------------- END ------------------
\end{document}